% Generated by scripts/rebuild_prompt_split.py.
\begin{xltabular}{\textwidth}{@{}>{\raggedright\arraybackslash}p{2.6cm}LL@{}}
\caption{The 30 further wordings of every task, by group: P is used to identify heads, Q only for the means patched in the test.}\label{tab:wording-split}\\
\toprule Group & P (identification) & Q (test) \\ \midrule \endfirsthead
\toprule Group & P (identification) & Q (test) \\ \midrule \endhead
\bottomrule \endlastfoot
\multicolumn{3}{@{}l}{\textit{antonym}} \\
cloze, names the task & \texttt{The opposite word for \{x\} is} \newline \texttt{Give the antonym of \{x\} as} \newline \texttt{For \{x\}, the reverse in meaning is} & \texttt{A word with the opposite meaning of \{x\} is} \newline \texttt{An antonym for \{x\} would be} \\
cloze, does not name it & \texttt{The contrasting word for \{x\} is} \newline \texttt{For \{x\}, the matching word with a contrary meaning is} & \texttt{The word that means the contrary of \{x\} is} \newline \texttt{A word whose meaning contrasts with \{x\} is} \newline \texttt{A single word that expresses a contrary idea to \{x\} is} \\
question--answer, names the task & \texttt{Provide the antonym for the given word. Q: \{x\}\nl A:} \newline \texttt{State the opposing word in meaning. Q: \{x\}\nl A:} & \texttt{Find the opposite meaning of this word. Q: \{x\}\nl A:} \newline \texttt{Give a reverse in meaning for the word. Q: \{x\}\nl A:} \newline \texttt{Identify an antonym of the supplied word. Q: \{x\}\nl A:} \\
question--answer, does not name it & \texttt{Respond with a single word expressing the contrasting meaning. Q: \{x\}\nl A:} \newline \texttt{Find a word whose meaning is contrary to the prompt word. Q: \{x\}\nl A:} \newline \texttt{Give a common term with the contrary sense. Q: \{x\}\nl A:} & \texttt{Provide a word with a contrary meaning to the given term. Q: \{x\}\nl A:} \newline \texttt{Write one word that means the contrary of the supplied word. Q: \{x\}\nl A:} \\
label, names the task & \texttt{Word: \{x\}\nl Antonym:} \newline \texttt{Given word: \{x\}\nl Opposing word:} \newline \texttt{Input word: \{x\}\nl Antonym output:} & \texttt{Term: \{x\}\nl Opposite:} \newline \texttt{Vocabulary item: \{x\}\nl Reverse meaning:} \\
label, does not name it & \texttt{Source word: \{x\}\nl Contrast word:} \newline \texttt{Prompt word: \{x\}\nl Contrary match:} & \texttt{Word: \{x\}\nl Contrary word:} \newline \texttt{Term: \{x\}\nl Contrasting term:} \newline \texttt{Given word: \{x\}\nl Contrary term:} \\
\midrule
\multicolumn{3}{@{}l}{\textit{synonym}} \\
cloze, names the task & \texttt{A similar word meaning \{x\} is} \newline \texttt{A term analogous in meaning to \{x\} is} & \texttt{A word equivalent to \{x\} is} \newline \texttt{A synonym for the word \{x\} is} \newline \texttt{An interchangeable word for \{x\} is} \\
cloze, does not name it & \texttt{One other word for \{x\} is} \newline \texttt{The word \{x\} can be restated as} \newline \texttt{To say \{x\} another way, use the word} & \texttt{Another word that means \{x\} is} \newline \texttt{For \{x\}, a different word with the same meaning is} \\
question--answer, names the task & \texttt{Give an equivalent word with the same meaning. Q: \{x\}\nl A:} \newline \texttt{Name an interchangeable word for the term. Q: \{x\}\nl A:} \newline \texttt{Respond with a word analogous in meaning. Q: \{x\}\nl A:} & \texttt{Provide one synonym for the given word. Q: \{x\}\nl A:} \newline \texttt{Find a similar word that means the same thing. Q: \{x\}\nl A:} \\
question--answer, does not name it & \texttt{Provide one different word that means the same thing. Q: \{x\}\nl A:} \newline \texttt{Return one other word that means what the given word means. Q: \{x\}\nl A:} & \texttt{Give another word with the same meaning. Q: \{x\}\nl A:} \newline \texttt{Respond with a single word that has the same definition. Q: \{x\}\nl A:} \newline \texttt{Write a different term carrying the same sense. Q: \{x\}\nl A:} \\
label, names the task & \texttt{Synonym prompt: \{x\}\nl Synonym:} \newline \texttt{Equivalent term: \{x\}\nl Equivalent:} & \texttt{Similar word: \{x\}\nl Similar:} \newline \texttt{Alike term: \{x\}\nl Alike word:} \newline \texttt{Analogous word: \{x\}\nl Analogous term:} \\
label, does not name it & \texttt{Term: \{x\}\nl Same sense:} \newline \texttt{Given word: \{x\}\nl Meaning match:} \newline \texttt{Entry: \{x\}\nl Same idea:} & \texttt{Word: \{x\}\nl Same meaning:} \newline \texttt{Source word: \{x\}\nl Other word:} \\
\midrule
\multicolumn{3}{@{}l}{\textit{singular-plural}} \\
cloze, names the task & \texttt{In the plural, \{x\} becomes} \newline \texttt{The plural version of \{x\} is} \newline \texttt{Write \{x\} as a plural noun} & \texttt{Change \{x\} to its plural form} \newline \texttt{For the noun \{x\}, the plural is} \\
cloze, does not name it & \texttt{To refer to several, write \{x\} as} \newline \texttt{For a count above one, the word \{x\} is written} & \texttt{The form of \{x\} for two or more is} \newline \texttt{When there is more than one, \{x\} becomes} \newline \texttt{If you mean two or more of them, \{x\} is} \\
question--answer, names the task & \texttt{Convert this singular English noun into its plural form. Q: \{x\}\nl A:} \newline \texttt{Give only the plural form of the noun shown. Q: \{x\}\nl A:} & \texttt{Provide the plural version of the given noun. Q: \{x\}\nl A:} \newline \texttt{Change the noun from singular to plural. Q: \{x\}\nl A:} \newline \texttt{Write the correct plural for this English noun. Q: \{x\}\nl A:} \\
question--answer, does not name it & \texttt{Write the noun as it would appear when referring to more than one item. Q: \{x\}\nl A:} \newline \texttt{Rewrite the word so it refers to more than one object. Q: \{x\}\nl A:} \newline \texttt{Change the noun from one of something to more than one. Q: \{x\}\nl A:} & \texttt{Give the noun form that means two or more of the same thing. Q: \{x\}\nl A:} \newline \texttt{Provide the standard English form used when the count is not one. Q: \{x\}\nl A:} \\
label, names the task & \texttt{Single Word: \{x\}\nl Plural Form:} \newline \texttt{Base Noun: \{x\}\nl Plural Version:} \newline \texttt{Singular Entry: \{x\}\nl Plural Entry:} & \texttt{Singular Noun: \{x\}\nl Plural Noun:} \newline \texttt{Noun Given: \{x\}\nl Plural:} \\
label, does not name it & \texttt{Term for one: \{x\}\nl Term for several:} \newline \texttt{Single noun: \{x\}\nl Several nouns:} & \texttt{Word for one: \{x\}\nl Word for many:} \newline \texttt{One noun: \{x\}\nl More than one noun:} \newline \texttt{One item: \{x\}\nl Many items:} \\
\midrule
\multicolumn{3}{@{}l}{\textit{present-past}} \\
cloze, names the task & \texttt{The simple past form of \{x\} is} \newline \texttt{Use \{x\} in the past tense as} & \texttt{In the past tense, \{x\} becomes} \newline \texttt{Converting \{x\} to past tense gives} \newline \texttt{When written for earlier time, \{x\} changes to} \\
cloze, does not name it & \texttt{If the action happened yesterday, \{x\} changes to} \newline \texttt{When referring to a finished action, the verb \{x\} becomes} \newline \texttt{The word \{x\} in the form used for a completed event is} & \texttt{For a sentence about something that happened before now, write \{x\} as} \newline \texttt{For talking about yesterday, the correct form of \{x\} is} \\
question--answer, names the task & \texttt{Convert the base verb to its past tense form. Q: \{x\}\nl A:} \newline \texttt{Give the simple past form of this English verb. Q: \{x\}\nl A:} \newline \texttt{Provide the verb form used for an earlier time. Q: \{x\}\nl A:} & \texttt{Change this verb from base form to past tense. Q: \{x\}\nl A:} \newline \texttt{State the past tense of the given base-form verb. Q: \{x\}\nl A:} \\
question--answer, does not name it & \texttt{Change the verb to the simple form used with yesterday. Q: \{x\}\nl A:} \newline \texttt{Provide the form of this verb for an event that happened before now. Q: \{x\}\nl A:} & \texttt{Write the English verb form for something done yesterday. Q: \{x\}\nl A:} \newline \texttt{Convert the base verb to the simple form used for a finished action. Q: \{x\}\nl A:} \newline \texttt{Supply the form this verb takes when reporting what someone did yesterday. Q: \{x\}\nl A:} \\
label, names the task & \texttt{Word: \{x\}\nl Past tense:} \newline \texttt{Action word: \{x\}\nl Past form:} & \texttt{Verb: \{x\}\nl Past tense:} \newline \texttt{Base verb: \{x\}\nl Past form:} \newline \texttt{Present verb: \{x\}\nl Past verb:} \\
label, does not name it & \texttt{Now Verb: \{x\}\nl Yesterday Verb:} \newline \texttt{Base Verb: \{x\}\nl Yesterday Form:} \newline \texttt{Dictionary Verb: \{x\}\nl Last Week Form:} & \texttt{Plain Verb: \{x\}\nl Yesterday Word:} \newline \texttt{Root Verb: \{x\}\nl Last Night Form:} \\
\midrule
\multicolumn{3}{@{}l}{\textit{country-capital}} \\
cloze, names the task & \texttt{The capital of \{x\} can be identified as} \newline \texttt{\{x\}'s capital city is} \newline \texttt{The city that serves as the capital of \{x\} is} & \texttt{For \{x\}, the capital city is} \newline \texttt{Name the capital city of \{x\} as} \\
cloze, does not name it & \texttt{The principal seat of government in \{x\} is} \newline \texttt{The place serving as \{x\}'s national administrative center is} & \texttt{The chief town of \{x\} is} \newline \texttt{The official chief town of \{x\} is} \newline \texttt{The official national seat for \{x\} is} \\
question--answer, names the task & \texttt{Identify the capital city for the country named. Q: \{x\}\nl A:} \newline \texttt{Name the country's capital city. Q: \{x\}\nl A:} & \texttt{Provide the capital of this country. Q: \{x\}\nl A:} \newline \texttt{Return only the capital city of the given country. Q: \{x\}\nl A:} \newline \texttt{What city is the capital of this country? Q: \{x\}\nl A:} \\
question--answer, does not name it & \texttt{Identify the primary seat of national government for this country. Q: \{x\}\nl A:} \newline \texttt{Provide the name of the country's official government seat. Q: \{x\}\nl A:} \newline \texttt{Name the place where this country's national government is based. Q: \{x\}\nl A:} & \texttt{Give the national seat of government for the named country. Q: \{x\}\nl A:} \newline \texttt{What is the official seat of government for this nation? Q: \{x\}\nl A:} \\
label, names the task & \texttt{Country: \{x\}\nl Capital city:} \newline \texttt{Country name: \{x\}\nl Capital name:} \newline \texttt{Sovereign country: \{x\}\nl Capital:} & \texttt{Nation: \{x\}\nl Capital:} \newline \texttt{Listed country: \{x\}\nl Capital city:} \\
label, does not name it & \texttt{Country: \{x\}\nl Seat of government:} \newline \texttt{Country name: \{x\}\nl Primary seat:} & \texttt{Nation: \{x\}\nl Government seat:} \newline \texttt{Sovereign state: \{x\}\nl Administrative center:} \newline \texttt{State: \{x\}\nl National seat:} \\
\midrule
\multicolumn{3}{@{}l}{\textit{country-currency}} \\
cloze, names the task & \texttt{For \{x\}, the official currency is} \newline \texttt{People in \{x\} use the currency} & \texttt{The currency used in \{x\} is} \newline \texttt{Name the currency of \{x\}} \newline \texttt{The money used in \{x\} is} \\
cloze, does not name it & \texttt{In \{x\}, people commonly pay with the} \newline \texttt{For purchases in \{x\}, the usual unit is the} \newline \texttt{The official tender recognized in \{x\} is the} & \texttt{The official unit of legal tender in \{x\} is the} \newline \texttt{The standard unit used for payments in \{x\} is the} \\
question--answer, names the task & \texttt{Identify the currency used in this country. Q: \{x\}\nl A:} \newline \texttt{Name the official currency for the country given. Q: \{x\}\nl A:} \newline \texttt{State the money used by the specified country. Q: \{x\}\nl A:} & \texttt{Provide the currency associated with this nation. Q: \{x\}\nl A:} \newline \texttt{Give the currency name for the following country. Q: \{x\}\nl A:} \\
question--answer, does not name it & \texttt{Identify the official legal tender used in the country Q: \{x\}\nl A:} \newline \texttt{Name the nationally accepted tender for everyday transactions Q: \{x\}\nl A:} & \texttt{Give the name of the national unit used for purchases there Q: \{x\}\nl A:} \newline \texttt{State the standard legal-tender unit for this country Q: \{x\}\nl A:} \newline \texttt{What official means of payment is used in this country? Q: \{x\}\nl A:} \\
label, names the task & \texttt{Nation: \{x\}\nl Currency:} \newline \texttt{Location: \{x\}\nl Money:} & \texttt{Country: \{x\}\nl Currency:} \newline \texttt{Country name: \{x\}\nl Official currency:} \newline \texttt{Country: \{x\}\nl Money used:} \\
label, does not name it & \texttt{Country: \{x\}\nl Legal tender:} \newline \texttt{Country name: \{x\}\nl Local tender:} \newline \texttt{Nation name: \{x\}\nl National tender:} & \texttt{Nation: \{x\}\nl Official tender:} \newline \texttt{State: \{x\}\nl Primary tender:} \\
\midrule
\multicolumn{3}{@{}l}{\textit{landmark-country}} \\
cloze, names the task & \texttt{The nation where \{x\} can be found is} \newline \texttt{\{x\} is located in the country of} \newline \texttt{Identify the country containing \{x\} as} & \texttt{For \{x\}, the country is} \newline \texttt{The country associated with the location of \{x\} is} \\
cloze, does not name it & \texttt{\{x\} is situated in the sovereign state of} \newline \texttt{At the level of sovereign states, \{x\} belongs to} & \texttt{The sovereign state where \{x\} is found is} \newline \texttt{The sovereign state containing \{x\} is} \newline \texttt{For \{x\}, the enclosing sovereign state is} \\
question--answer, names the task & \texttt{Identify the country for this place Q: \{x\}\nl A:} \newline \texttt{Name the country associated with this location Q: \{x\}\nl A:} & \texttt{State the nation where this site is located Q: \{x\}\nl A:} \newline \texttt{Give the country containing this landmark Q: \{x\}\nl A:} \newline \texttt{Determine the country of this building or institution Q: \{x\}\nl A:} \\
question--answer, does not name it & \texttt{In what sovereign state is this place situated? Q: \{x\}\nl A:} \newline \texttt{Name the sovereign state where this institution is based. Q: \{x\}\nl A:} \newline \texttt{Which sovereign state is home to this named location? Q: \{x\}\nl A:} & \texttt{Identify the sovereign state that contains this landmark. Q: \{x\}\nl A:} \newline \texttt{Give the sovereign state where this building or site stands. Q: \{x\}\nl A:} \\
label, names the task & \texttt{Place nation: \{x\}\nl Nation:} \newline \texttt{Institution country: \{x\}\nl Country:} \newline \texttt{Site country lookup: \{x\}\nl Country:} & \texttt{Landmark country: \{x\}\nl Country:} \newline \texttt{Building nation: \{x\}\nl Nation:} \\
label, does not name it & \texttt{Site: \{x\}\nl Sovereign state:} \newline \texttt{Location: \{x\}\nl Political state:} & \texttt{Landmark: \{x\}\nl Recognized state:} \newline \texttt{Place name: \{x\}\nl State:} \newline \texttt{Entry: \{x\}\nl Sovereign realm:} \\
\midrule
\multicolumn{3}{@{}l}{\textit{park-country}} \\
cloze, names the task & \texttt{The country associated with \{x\} is} \newline \texttt{The country that includes \{x\} is} & \texttt{For \{x\}, the country is} \newline \texttt{The country containing \{x\} is} \newline \texttt{The country home to \{x\} is} \\
cloze, does not name it & \texttt{The sovereign state where \{x\} lies is} \newline \texttt{The sovereign territory that includes \{x\} is} \newline \texttt{The state whose territory includes \{x\} is} & \texttt{The political state containing \{x\} is} \newline \texttt{\{x\} is found within the borders of} \\
question--answer, names the task & \texttt{Provide the country where the named park or reserve can be found. Q: \{x\}\nl A:} \newline \texttt{Determine the country that contains this park or nature reserve. Q: \{x\}\nl A:} \newline \texttt{State the country associated with the following park or nature reserve. Q: \{x\}\nl A:} & \texttt{Identify the country for this park or nature reserve. Q: \{x\}\nl A:} \newline \texttt{Name the country in which this protected natural area is located. Q: \{x\}\nl A:} \\
question--answer, does not name it & \texttt{Provide the independent state that includes this protected area Q: \{x\}\nl A:} \newline \texttt{Give the independent political state for the listed protected site Q: \{x\}\nl A:} & \texttt{Determine the sovereign state where this park or reserve is found Q: \{x\}\nl A:} \newline \texttt{State the sovereign jurisdiction containing this park or nature reserve Q: \{x\}\nl A:} \newline \texttt{Which sovereign state is home to this named park or reserve Q: \{x\}\nl A:} \\
label, names the task & \texttt{Protected area: \{x\}\nl Country:} \newline \texttt{Park name: \{x\}\nl Country:} & \texttt{Park or reserve: \{x\}\nl Country:} \newline \texttt{Nature site: \{x\}\nl Country:} \newline \texttt{Reserve name: \{x\}\nl Country:} \\
label, does not name it & \texttt{Protected area: \{x\}\nl Sovereign state:} \newline \texttt{Reserve name: \{x\}\nl State:} \newline \texttt{Place name: \{x\}\nl Sovereign location:} & \texttt{Park name: \{x\}\nl Located in:} \newline \texttt{Natural site: \{x\}\nl Jurisdiction:} \\
\midrule
\multicolumn{3}{@{}l}{\textit{product-company}} \\
cloze, names the task & \texttt{The manufacturer of \{x\} is} \newline \texttt{The corporation that makes \{x\} is} \newline \texttt{For \{x\}, the manufacturer is} & \texttt{The company responsible for making \{x\} is} \newline \texttt{The firm behind \{x\} is} \\
cloze, does not name it & \texttt{The maker of \{x\} is} \newline \texttt{\{x\} is produced by} & \texttt{\{x\} is made by} \newline \texttt{The business that makes \{x\} is} \newline \texttt{The producer of \{x\} is} \\
question--answer, names the task & \texttt{State which firm produces this product. Q: \{x\}\nl A:} \newline \texttt{Find the company associated with making this product. Q: \{x\}\nl A:} & \texttt{Identify the company that makes the named product. Q: \{x\}\nl A:} \newline \texttt{Name the manufacturer of the specified product. Q: \{x\}\nl A:} \newline \texttt{Give the corporation behind this product. Q: \{x\}\nl A:} \\
question--answer, does not name it & \texttt{Identify who makes the named product. Q: \{x\}\nl A:} \newline \texttt{Name the maker of this product. Q: \{x\}\nl A:} \newline \texttt{Which business produces this product? Q: \{x\}\nl A:} & \texttt{State who creates the named item. Q: \{x\}\nl A:} \newline \texttt{Tell me the maker behind this item. Q: \{x\}\nl A:} \\
label, names the task & \texttt{Product: \{x\}\nl Company:} \newline \texttt{Item: \{x\}\nl Manufacturer:} \newline \texttt{Named product: \{x\}\nl Corporation:} & \texttt{Product name: \{x\}\nl Firm:} \newline \texttt{Listed item: \{x\}\nl Making company:} \\
label, does not name it & \texttt{Item name: \{x\}\nl Creator:} \newline \texttt{Goods: \{x\}\nl Source:} & \texttt{Product: \{x\}\nl Maker:} \newline \texttt{Item: \{x\}\nl Made by:} \newline \texttt{Product name: \{x\}\nl Producer:} \\
\midrule
\multicolumn{3}{@{}l}{\textit{next-item}} \\
cloze, names the task & \texttt{In the sequence, the next entry after \{x\} is} \newline \texttt{The following item after \{x\} would be} & \texttt{The item immediately following \{x\} is} \newline \texttt{Right after \{x\} comes} \newline \texttt{The next sequence item following \{x\} is} \\
cloze, does not name it & \texttt{Advancing \{x\} by one place gives} \newline \texttt{The immediate successor of \{x\} is} \newline \texttt{Move one step forward from \{x\} to} & \texttt{The item that succeeds \{x\} is} \newline \texttt{One place later than \{x\} comes} \\
question--answer, names the task & \texttt{Identify the next item in its familiar sequence. Q: \{x\}\nl A:} \newline \texttt{Give the item that comes immediately after this one. Q: \{x\}\nl A:} \newline \texttt{State the next entry in the appropriate known ordering. Q: \{x\}\nl A:} & \texttt{Provide the following item in the relevant standard sequence. Q: \{x\}\nl A:} \newline \texttt{For this sequence item, answer with what follows after it. Q: \{x\}\nl A:} \\
question--answer, does not name it & \texttt{Write the item one position later in its usual series Q: \{x\}\nl A:} \newline \texttt{State the item that comes one place later in the customary progression Q: \{x\}\nl A:} & \texttt{Provide the immediate successor in the appropriate familiar order Q: \{x\}\nl A:} \newline \texttt{Identify the entry reached by advancing one step in the standard sequence Q: \{x\}\nl A:} \newline \texttt{Return the direct successor from the relevant ordered set Q: \{x\}\nl A:} \\
label, names the task & \texttt{Sequence entry: \{x\}\nl Following entry:} \newline \texttt{Given item: \{x\}\nl Item after:} & \texttt{Current item: \{x\}\nl Next item:} \newline \texttt{Present value: \{x\}\nl Next value:} \newline \texttt{Starting point: \{x\}\nl Following point:} \\
label, does not name it & \texttt{Given item: \{x\}\nl Immediate successor:} \newline \texttt{Sequence entry: \{x\}\nl Succeeding entry:} \newline \texttt{Present term: \{x\}\nl Subsequent term:} & \texttt{Current: \{x\}\nl Successor:} \newline \texttt{Starting item: \{x\}\nl Successor item:} \\
\midrule
\multicolumn{3}{@{}l}{\textit{prev-item}} \\
cloze, names the task & \texttt{The item immediately before \{x\} is} \newline \texttt{In the relevant familiar sequence, the previous item to \{x\} is} \newline \texttt{The last item before \{x\} in its sequence is} & \texttt{What comes right before \{x\} is} \newline \texttt{The earlier item in the sequence for \{x\} is} \\
cloze, does not name it & \texttt{In its usual ordered set, the item one step back from \{x\} is} \newline \texttt{Step back once from \{x\} in the usual order to get} & \texttt{Move one place backward from \{x\} and the result is} \newline \texttt{For the standard sequence containing \{x\}, the neighbor immediately to its left is} \newline \texttt{The entry that comes just to the left of \{x\} is} \\
question--answer, names the task & \texttt{Return only the previous item in the familiar ordering for the supplied entry. Q: \{x\}\nl A:} \newline \texttt{For a number, letter, Roman numeral, day, or month, give the previous item. Q: \{x\}\nl A:} & \texttt{Identify the sequence member immediately before the given number, letter, Roman numeral, weekday, or month. Q: \{x\}\nl A:} \newline \texttt{Find what comes earlier by one step in the relevant sequence. Q: \{x\}\nl A:} \newline \texttt{State the item immediately before the provided sequence item. Q: \{x\}\nl A:} \\
question--answer, does not name it & \texttt{Return the item one step to the left in the normal sequence. Q: \{x\}\nl A:} \newline \texttt{Report the item directly to the left in the recognized sequence. Q: \{x\}\nl A:} \newline \texttt{Move one place backward in the appropriate sequence and output that entry. Q: \{x\}\nl A:} & \texttt{For the given token, write the entry directly to its left in the usual order. Q: \{x\}\nl A:} \newline \texttt{Supply the item that is one position lower in the applicable sequence. Q: \{x\}\nl A:} \\
label, names the task & \texttt{Current item: \{x\}\nl Previous item:} \newline \texttt{Given element: \{x\}\nl Earlier element:} \newline \texttt{Present value: \{x\}\nl Previous value:} & \texttt{Sequence entry: \{x\}\nl Comes before:} \newline \texttt{Reference item: \{x\}\nl Item before:} \\
label, does not name it & \texttt{Current symbol: \{x\}\nl Move back one:} \newline \texttt{Current term: \{x\}\nl Immediate left:} & \texttt{Entry: \{x\}\nl One step back:} \newline \texttt{Item: \{x\}\nl Step back:} \newline \texttt{Sequence entry: \{x\}\nl Left neighbor:} \\
\midrule
\multicolumn{3}{@{}l}{\textit{english-french}} \\
cloze, names the task & \texttt{In French, \{x\} is} \newline \texttt{Please translate \{x\} into French as} & \texttt{The word \{x\} in French is} \newline \texttt{The French equivalent of \{x\} is} \newline \texttt{For the English word \{x\}, the French translation is} \\
cloze, does not name it & \texttt{In the language of Paris, \{x\} is} \newline \texttt{In the official language of Paris, \{x\} becomes} \newline \texttt{The Parisian term matching \{x\} is} & \texttt{The word for \{x\} in the language of Molière is} \newline \texttt{A Paris native would say \{x\} as} \\
question--answer, names the task & \texttt{Write the English word in French. Q: \{x\}\nl A:} \newline \texttt{Convert this English word to French. Q: \{x\}\nl A:} \newline \texttt{State the French equivalent of the word. Q: \{x\}\nl A:} & \texttt{Provide the French translation of the given English word. Q: \{x\}\nl A:} \newline \texttt{Give the corresponding French word. Q: \{x\}\nl A:} \\
question--answer, does not name it & \texttt{Give the equivalent of this English word in the main language spoken in Paris. Q: \{x\}\nl A:} \newline \texttt{Write the usual word in the language of Molière for this English term. Q: \{x\}\nl A:} & \texttt{Supply the standard Parisian-language counterpart for this English word. Q: \{x\}\nl A:} \newline \texttt{Convert the English word into the official tongue used in Paris. Q: \{x\}\nl A:} \newline \texttt{Provide the matching word in the national language heard across Paris. Q: \{x\}\nl A:} \\
label, names the task & \texttt{English to French: \{x\}\nl French:} \newline \texttt{French translation: \{x\}\nl Translation:} & \texttt{English word: \{x\}\nl French equivalent:} \newline \texttt{Source word: \{x\}\nl French word:} \newline \texttt{Input English: \{x\}\nl French output:} \\
label, does not name it & \texttt{English word: \{x\}\nl FR equivalent:} \newline \texttt{English entry: \{x\}\nl FR entry:} \newline \texttt{EN item: \{x\}\nl FR item:} & \texttt{EN term: \{x\}\nl FR term:} \newline \texttt{Source word: \{x\}\nl FR word:} \\
\midrule
\multicolumn{3}{@{}l}{\textit{english-german}} \\
cloze, names the task & \texttt{In German, \{x\} means} \newline \texttt{The English word \{x\} translates into German as} \newline \texttt{A German rendering of \{x\} is} & \texttt{The Deutsch equivalent of \{x\} is} \newline \texttt{For the word \{x\}, the German translation is} \\
cloze, does not name it & \texttt{For the IETF language tag de-DE, the English word \{x\} is} \newline \texttt{The de-DE equivalent of the English word \{x\} is} & \texttt{In the language identified by ISO 639-1 code de, the English word \{x\} is} \newline \texttt{Using ISO code de for the target language, the English word \{x\} is} \newline \texttt{The English word \{x\} in the language with two-letter code de is} \\
question--answer, names the task & \texttt{Provide the German translation for the English word Q: \{x\}\nl A:} \newline \texttt{Give the German equivalent of this English word Q: \{x\}\nl A:} & \texttt{Translate this English word into German Q: \{x\}\nl A:} \newline \texttt{Write the German word that matches the English term Q: \{x\}\nl A:} \newline \texttt{Convert the English word to its German translation Q: \{x\}\nl A:} \\
question--answer, does not name it & \texttt{Provide the standard Berlin-language equivalent of this English word Q: \{x\}\nl A:} \newline \texttt{Give the word a Berlin speaker would use for this English word Q: \{x\}\nl A:} \newline \texttt{Supply the Berlin standard counterpart for the given English word Q: \{x\}\nl A:} & \texttt{State the common form in the language of Berlin for the English word Q: \{x\}\nl A:} \newline \texttt{Write the matching term in the official language used in Berlin for this English word Q: \{x\}\nl A:} \\
label, names the task & \texttt{English Word: \{x\}\nl German Translation:} \newline \texttt{Input Word: \{x\}\nl German:} \newline \texttt{Word: \{x\}\nl Deutsch:} & \texttt{Source Word: \{x\}\nl German Equivalent:} \newline \texttt{English Term: \{x\}\nl German Word:} \\
label, does not name it & \texttt{Input word: \{x\}\nl DE output:} \newline \texttt{Source lexeme: \{x\}\nl DE lexeme:} & \texttt{English word: \{x\}\nl DE word:} \newline \texttt{Source term: \{x\}\nl DE equivalent:} \newline \texttt{English entry: \{x\}\nl DE entry:} \\
\midrule
\multicolumn{3}{@{}l}{\textit{english-spanish}} \\
cloze, names the task & \texttt{\{x\} rendered in Spanish is} \newline \texttt{A Spanish translation for \{x\} is} & \texttt{In Spanish, the word \{x\} is} \newline \texttt{The Spanish equivalent of \{x\} is} \newline \texttt{The word \{x\} translates into Spanish as} \\
cloze, does not name it & \texttt{In Castilian, the English word \{x\} is} \newline \texttt{The Castilian equivalent of \{x\} is} \newline \texttt{For the English word \{x\}, the Castilian equivalent is} & \texttt{The Castilian word for \{x\} is} \newline \texttt{\{x\} in Castilian is} \\
question--answer, names the task & \texttt{Provide the Spanish equivalent of the English word. Q: \{x\}\nl A:} \newline \texttt{What is the Spanish translation for this English word? Q: \{x\}\nl A:} \newline \texttt{Convert the following English vocabulary item to Spanish. Q: \{x\}\nl A:} & \texttt{Translate this English word into Spanish. Q: \{x\}\nl A:} \newline \texttt{Write the given English word in Spanish. Q: \{x\}\nl A:} \\
question--answer, does not name it & \texttt{Give the Castilian equivalent of the English word. Q: \{x\}\nl A:} \newline \texttt{Supply the Castilian counterpart for this English word. Q: \{x\}\nl A:} & \texttt{Provide the matching Castilian word for the given English term. Q: \{x\}\nl A:} \newline \texttt{Write the Castilian word that has the same meaning. Q: \{x\}\nl A:} \newline \texttt{Answer with the equivalent term in Castilian. Q: \{x\}\nl A:} \\
label, names the task & \texttt{Word To Translate: \{x\}\nl Spanish:} \newline \texttt{Original Word: \{x\}\nl Spanish Word:} & \texttt{English Word: \{x\}\nl Spanish Translation:} \newline \texttt{Source Term: \{x\}\nl Spanish Equivalent:} \newline \texttt{English Vocabulary: \{x\}\nl Spanish Vocabulary:} \\
label, does not name it & \texttt{English word: \{x\}\nl Castilian word:} \newline \texttt{English term: \{x\}\nl Castilian equivalent:} \newline \texttt{English vocabulary: \{x\}\nl Castilian vocabulary:} & \texttt{English item: \{x\}\nl Castilian counterpart:} \newline \texttt{English entry: \{x\}\nl Castilian entry:} \\

\end{xltabular}
